\section{Full permanent multiplicities from detected Taylor coefficients}

\subsection*{Finding F-021: the internal-padding permanent jet map}
Let \(r\geq2\), \(R=r-1\), \(\ell\geq0\), and set
\[
 b=2\ell+2,\qquad h=\lceil\sqrt R\rceil,\qquad
 J_b(r)=2h\binom{R+\ell}{\ell}.
\]
Suppose \(n\geq m\geq J_b(r)\). For every outer degree
\(d\) and every partition \(\lambda\vdash nd\) with
\(\ell(\lambda)\leq r\) and \(\lambda_2\leq b\),
\begin{equation}\label{eq:F021}
 q_\lambda(m,n)=K_n(\lambda)=a_\lambda(d[n]).
\end{equation}
Here \(q\) uses the internal padding
\(X_{11}^{n-m}\operatorname{per}_m\), exactly as in the
supplied IP and BIP sources. The statement includes every
highest-weight copy, and imposes no upper bound on \(d\).
The usual setting \(n>m\) is a special case. This is a
variant derived in this run from M-013, without a priority claim.

\begin{proof}
Use the two separate families of rooted monomial trees
and the bridges of M-013. Thus \(B(z)\) is the adjacency
matrix of an acyclic directed graph on \(J=J_b(r)\)
vertices. Its inward-to-outward path from \((\alpha,a)\)
to \((\beta,c)\), when the bridge labeled \(z_{ac}\)
is present, is unique and has product
\[
                         z^\alpha z_{ac}z^\beta.
                                                               \tag{F-021.1}
\]
All nonempty paths have positive labels. Put
\(A_0(x_1,z)=x_1 I+B(z)\), and append an \(x_1 I\)
block if \(m>J\). Choose any one vertex as the distinguished
first vertex, so \((A_0)_{11}=x_1\).
Since the graph is acyclic, the only permutation contributing
to the permanent of \(I+B\) is its diagonal permutation.
Consequently \(\operatorname{per}A_0=x_1^m\).

The permanent cofactor obtained by deleting row \(i\)
and column \(j\) from \(I+B\) has a useful exact
description. If \(i=j\), it equals one. If \(i\ne j\),
it is the sum of the products of all directed paths from
\(j\) to \(i\). Indeed a term in this minor chooses one
outgoing edge at each vertex other than \(i\), and one
incoming edge at each vertex other than \(j\). Its graph
therefore consists of a directed path from \(j\) to \(i\),
loops, and possibly directed cycles. Acyclicity excludes
all nontrivial cycles, and every such path together with
loops on its unused vertices gives exactly one minor term.
All signs are positive. In particular, the cofactors on
the chosen inward-to-outward pairs are the single monomials
in (F-021.1).

Every monomial of degree between one and \(2\ell+1\)
has the form (F-021.1): remove one variable for a bridge
and split the remaining multiset into two parts of degrees
at most \(\ell\). Diagonal cofactors give the constant
monomial. The linear-entry differential
\[
 d(\operatorname{per})_{A_0}(H E_{ij})
       =H(1,z)\operatorname{per}\bigl((I+B)_{\widehat i,
                                                        \widehat j}\bigr)
                                                               \tag{F-021.2}
\]
therefore spans every monomial through degree \(b\), by
taking \(H=x_1\) through degree \(b-1\) and multiplying
a degree-\(b-1\) cofactor by a variable for the final degree.

The distinguished entry must remain exactly \(x_1\),
so none of its coefficients is varied. This does not remove
any required image direction. One has \(J\geq b\geq2\),
and hence there is another diagonal entry to supply the
constant monomial, as well as any linear monomial if needed.
All positive-degree path cofactors used above come from
off-diagonal entries. Distinct monomial image directions
can be chosen from distinct elementary entry-coefficient
coordinates, each with coefficient one. Thus the map
\[
 \{A:\ A_{11}=x_1\}\longrightarrow
 \mathbb Z[z_1,\ldots,z_R]_{\leq b},\qquad
 A\longmapsto\operatorname{per}A(1,z)\pmod{(z)^{b+1}}
                                                               \tag{F-021.3}
\]
has an identity Jacobian minor at \(A_0\), over
\(\mathbb Z\). Restricting to those source coordinates
and translating origins gives \(F(t)=t+O(t^2)\).
The lowest nonzero homogeneous part of a polynomial survives
substitution into \(F\). Its polynomial pullback is therefore
injective, after every commutative base change. In characteristic
zero this in particular proves dominance of (F-021.3).

For the internally padded permanent, substitution of the
pencil \(A\) with \(A_{11}=x_1\) gives
\[
                         x_1^{n-m}\operatorname{per}A(x_1,z).
\]
At \(x_1=1\) its detected coefficients are exactly those
of (F-021.3). All these forms belong to the permanent orbit
closure: a linear substitution into the matrix entries is
an endomorphism substitution, and invertible endomorphisms
are dense. The variable space fits because
\(m\geq2\lceil\sqrt{r-1}\rceil\) implies \(r\leq m^2\).

Now let \(w\) be any highest-weight vector of weight
\(\lambda\) in the permanent ideal. It uses only the
first \(r\) coordinates. F-016 with \(j=2\) says that
every coefficient generator it uses has degree at most
\(\lambda_2\leq b\) outside the first coordinate.
Thus \(w\) is a polynomial in the detected jet coordinates.
It vanishes on the internally padded pencil family, so
the proved pullback injection forces \(w=0\).
The entire ideal multiplicity space vanishes, and hence
\(q_\lambda=a_\lambda\). M-013 applies as well because
\(n\geq m\geq J_b(r)\), giving \(K_n=a_\lambda\).
This proves \eqref{eq:F021}.
\end{proof}

\paragraph{An additional order-five case.}
The cyclic-star construction of M-015 also has a permanent
version. For \(r\geq4\) and \(n\geq m\geq r^2\),
\begin{equation}\label{eq:F021five}
 \ell(\lambda)\leq r,\quad\lambda_2\leq5
       \quad\Longrightarrow\quad q_\lambda=K_n(\lambda)=a_\lambda.
\end{equation}
To verify this variant, keep the same central vertex, core
vertices, leaves, and positive edge labels as M-015, and
use \(A=x_1I+B\). Every nontrivial directed cycle is a
two-edge cycle through the central vertex. Thus at \(x_1=1\)
the base permanent is
\(D_+=1+\sum_i z_i z_{\sigma(i)}\).

The cofactor for the central diagonal is one. A cofactor
whose open path runs from the central vertex to a core
vertex is the corresponding single variable. Between distinct
core vertices \(v_i,v_k\), the only open path is
\(v_i\longrightarrow0\longrightarrow v_k\), with product
\(z_i z_{\sigma(k)}\). The chosen inward-to-outward leaf
cofactors have \(i\ne\sigma^{-1}(b)\), and their only
open path is
\[
 p_{ij}\longrightarrow v_i\longrightarrow0
         \longrightarrow v_{\sigma^{-1}(b)}\longrightarrow q_{ab},
\]
with product \(z_i z_j z_a z_b\). In each positive-degree
case the path uses the central vertex, so no disjoint
nontrivial cycle can contribute. All unused vertices take
their diagonal loops. These are therefore pure permanent
cofactors, with coefficient one. The complete degree-two
and degree-four coverage proved in M-015 applies unchanged.
Multiplication by a linear entry perturbation supplies the
missing degrees three and five.

Choose an inward leaf as the distinguished vertex and hold
its diagonal equal to \(x_1\). None of the selected directions
varies that diagonal: the constant comes from the central
diagonal and all the other selected cofactors are off diagonal.
The same integral unit-Jacobian and internal-padding argument
as above proves \(q_\lambda=a_\lambda\). M-015 gives
\(K_n=a_\lambda\), proving \eqref{eq:F021five}.

\paragraph{A strict improvement over every exact two-tree support cover.}
For the full simplex
\(\Gamma=\{\alpha\in\mathbb Z_{\geq0}^R:|\alpha|\leq4\}\),
\(R\geq4\), the finite cover invariant of F-018 satisfies
\begin{equation}\label{eq:F021cover}
             \kappa(\Gamma)\geq
                          1+\left\lceil\frac{R(R-1)}{12}\right\rceil.
\end{equation}
To prove this, let \(P,Q\) be any admissible downward-closed
sets. For a squarefree degree-four exponent \(\gamma\),
write \(\gamma=\alpha+\beta+\delta\), where
\(\alpha\in P\), \(\beta\in Q\), and
\(\delta\in\{0,e_1,\ldots,e_R\}\). At least one of
\(\alpha,\beta\) has degree at least two, and its
downward-closed set contains a squarefree pair supported
inside \(\gamma\). Let \(E\) be the set of pairs in
\(P\cup Q\). Every four-subset contains a member of \(E\),
so double-counting gives
\[
 |E|\binom{R-2}{2}\geq\binom R4,
          \qquad |E|\geq\left\lceil R(R-1)/12\right\rceil.
\]
Both \(P,Q\) contain zero; consequently
\(|P|+|Q|-1\geq|P\cup Q|\geq1+|E|\), proving the claim.

Take \(r=1024\), \(R=1023\), and \(b=4\). Then
\[
 J_4(1024)=2\cdot32\cdot1024=65536,
                    \qquad \kappa(\Gamma)\geq87127.
\]
With \(m=65536\) and \(n=65537\), F-021 applies to
every type with second row at most four, whereas neither
\(m\) nor \(n\) reaches the exact-cover lower bound for
this simplex. This is a limitation of the F-018 architecture,
not a lower bound for arbitrary permanent or determinant
representations.

This strict threshold comparison contains positive higher
multiplicity examples. Set \(d=2r\) and
\[
                  \lambda=(nd-4(r-1),4^{r-1}).
                                                               \tag{F-021.4}
\]
Its full nested support at either base degree four or any
larger degree is precisely \(\Gamma\). Its ambient
multiplicity is at least two, as follows. Let \(c_0(F)=F(e_1)\)
and let \(W_r(F)\) be the leading \(r\)-by-\(r\)
Hessian determinant at \(e_1\), as in O-011.
The polynomial \(W_r^2\) has degree \(2r\) and weight
\(\lambda\). It is not divisible by \(c_0\): the independent
coefficient variables can prescribe a Hessian with its
\((1,1)\) entry zero, its \((1,2),(2,1)\) entries one,
its \((2,2)\) entry zero, and an identity block in the
remaining coordinates. At this point \(c_0=0\) but
\(W_r=-1\ne0\), after harmless nonzero coefficient scalings.

F-011, with outer degree \(r\) and \(k=4\), supplies
a nonzero highest-weight polynomial \(I\) of degree \(r\)
and weight \((r(n-4)+4,4^{r-1})\). Therefore
\(I c_0^r\) also has degree \(2r\) and weight \(\lambda\).
It is nonzero and divisible by \(c_0\), whereas \(W_r^2\)
is not. These two highest-weight polynomials are linearly
independent. Hence (F-021.4) at the stated parameters has
\[
                         2\leq q_\lambda=K_n(\lambda)=a_\lambda.
\]
No numerical value larger than this lower bound is asserted.

\paragraph{Grothendieck dictionary and remaining scope.}
The permanent minor is a positive sum of path monomials,
and the resulting coefficient maps are integral morphisms
of affine schemes. The truncated polynomial module is finite
locally free. Its coordinate algebra injects into the source
parameter algebra after arbitrary base change by the unit
Jacobian argument. Schur highest weights and the characteristic-zero
quotient decomposition identify the complete multiplicity spaces;
these are the comodule and Schur-functor operations of SGA~3
I4.7.2--4.7.3 and SGA~6 VI3.3. This geometric injection
contains information not supplied by an equality of virtual
\(K_0\)-classes. The conclusion determines additional full
permanent multiplicities; it does not assert a comparison
outside the stated jet-size condition or solve the remaining
general evaluation-rank problem.
